\chapter{Chapter05 — Unate Functions}
%%% generated by the blueprint pipeline from FABL.Chapter05.UnateFunctions
\section{Overview}

Declarations of the Lean module \texttt{FABL.Chapter05.UnateFunctions}.

\section{Declarations}

\begin{definition}[Input reorientation equivalence]
\lean{FABL.reorientInputEquiv}
\statementsource{Analysis-of-Boolean-Functions-by-Ryan-ODonnell}{pdf-ffbe9c7b3306-p039-b005}
\label{FABL.reorientInputEquiv}
\leanok
For a point $\sigma$ of the $n$-dimensional Boolean cube, reorient the cube by sending
each point $x$ to the point whose $i$th coordinate is $\sigma_i x_i$. This map is a
bijection, and applying it twice returns the original point.
\end{definition}

\begin{theorem}[Action of input reorientation]
\lean{FABL.reorientInputEquiv_apply}
\label{FABL.reorientInputEquiv_apply}
\leanok
\uses{FABL.reorientInputEquiv}
\difficulty{1}
Let $\sigma$ and $x$ be sign vectors in $\{\pm1\}^n$. Applying the reorientation
determined by $\sigma$ to $x$ via the associated equivalence is the same as applying
that reorientation directly to $x$.
\end{theorem}

\begin{theorem}[Reorientation commutes with coordinate flipping]
\lean{FABL.reorientInput_flipCoordinate}
\label{FABL.reorientInput_flipCoordinate}
\leanok
\uses{FABL.flipCoordinate, FABL.setCoordinate}
\difficulty{4}
Let $\sigma,x\in\{\pm1\}^n$, and let $i$ be a coordinate. Reorienting the input after
flipping coordinate $i$ gives the same sign vector as flipping coordinate $i$ after
reorientation.
\end{theorem}

\begin{theorem}[Influence under input reorientation]
\lean{FABL.booleanInfluence_comp_reorientInput}
\label{FABL.booleanInfluence_comp_reorientInput}
\leanok
\uses{FABL.BooleanFunction, FABL.IsPivotal, FABL.booleanInfluence, FABL.flipCoordinate, FABL.reorientInputEquiv, FABL.reorientInput_flipCoordinate, FABL.uniformProbability}
\difficulty{4}
Let $n$ be a nonnegative integer, let $f$ be a function from the $n$-dimensional sign
cube to $\{\pm1\}$, let $\sigma$ be a sign vector in that cube, and let $i$ be a
coordinate. The influence of $f$ in coordinate $i$ is unchanged when the input is
reoriented by $\sigma$.
\end{theorem}

\begin{definition}[Coordinatewise monotonicity]
\lean{FABL.IsMonotoneInCoordinate}
\statementsource{Analysis-of-Boolean-Functions-by-Ryan-ODonnell}{
  pdf-ffbe9c7b3306-p045-b002,
  pdf-ffbe9c7b3306-p047-b006
}
\label{FABL.IsMonotoneInCoordinate}
\leanok
\uses{FABL.BooleanFunction, FABL.setCoordinate}
Let $n$ be a nonnegative integer, let $f$ be a Boolean function on the $n$-dimensional
cube, and choose a coordinate $i$. The function is monotone in coordinate $i$ if, for
every input $x$, its value when coordinate $i$ is set to $-1$ is at most its value when
coordinate $i$ is set to $1$.
\end{definition}

\begin{definition}[Coordinatewise antimonotonicity]
\lean{FABL.IsAntimonotoneInCoordinate}
\statementsource{Analysis-of-Boolean-Functions-by-Ryan-ODonnell}{pdf-ffbe9c7b3306-p061-b008}
\label{FABL.IsAntimonotoneInCoordinate}
\leanok
\uses{FABL.BooleanFunction, FABL.setCoordinate}
Let $n$ be a nonnegative integer, let $f$ be a Boolean function on the $n$-dimensional
cube, and choose a coordinate $i$. The function is antimonotone in coordinate $i$ if,
for every input $x$, its value when coordinate $i$ is set to $1$ is at most its value
when coordinate $i$ is set to $-1$.
\end{definition}

\begin{definition}[Coordinatewise unateness]
\lean{FABL.IsUnateInCoordinate}
\statementsource{Analysis-of-Boolean-Functions-by-Ryan-ODonnell}{pdf-ffbe9c7b3306-p061-b008}
\label{FABL.IsUnateInCoordinate}
\leanok
\uses{FABL.BooleanFunction, FABL.IsAntimonotoneInCoordinate, FABL.IsMonotoneInCoordinate}
Let $n$ be a nonnegative integer, let $f$ be a Boolean function on the $n$-dimensional
cube, and choose a coordinate $i$. The function is unate in coordinate $i$ if it is
either monotone or antimonotone in that coordinate.
\end{definition}

\begin{definition}[Unate Boolean function]
\lean{FABL.IsUnate}
\statementsource{Analysis-of-Boolean-Functions-by-Ryan-ODonnell}{
  pdf-ffbe9c7b3306-p061-b008,
  pdf-ffbe9c7b3306-p132-b002
}
\label{FABL.IsUnate}
\leanok
\uses{FABL.BooleanFunction, FABL.IsUnateInCoordinate}
A Boolean function on $n$ coordinates is unate if it is unate in each coordinate.
\end{definition}

\begin{theorem}[Absolute derivative equals squared derivative]
\lean{FABL.abs_discreteDerivative_toReal_eq_sq}
\proofsource{Analysis-of-Boolean-Functions-by-Ryan-ODonnell}{pdf-ffbe9c7b3306-p047-b003}
\label{FABL.abs_discreteDerivative_toReal_eq_sq}
\leanok
\uses{FABL.BooleanFunction, FABL.BooleanFunction.toReal, FABL.discreteDerivative, FABL.discreteDerivative_apply, FABL.setCoordinate}
\difficulty{4}
Let $f$ be a Boolean function on the sign cube $\\{−1,1\\}^n$, let $i$ be a coordinate,
and let $x$ be a point of the cube. The absolute value of the discrete derivative of the
real-valued version of $f$ at $x$ in coordinate $i$ equals the square of that
derivative.
\end{theorem}

\begin{theorem}[Singleton Fourier coefficient and influence]
\lean{FABL.abs_fourierCoeff_singleton_le_influence}
\label{FABL.abs_fourierCoeff_singleton_le_influence}
\leanok
\uses{FABL.BooleanFunction, FABL.BooleanFunction.toReal, FABL.abs_discreteDerivative_toReal_eq_sq, FABL.discreteDerivative, FABL.fourierCoeff, FABL.influence, FABL.mean, FABL.mean_discreteDerivative_eq_fourierCoeff_singleton}
\difficulty{4}
For any Boolean function $f$ on the sign cube $\\{-1,1\\}^n$ and any coordinate $i$, the
absolute value of its Fourier coefficient at $\\{i\\}$ is at most its influence at $i$.
\end{theorem}

\begin{theorem}[Nonnegativity of the derivative square minus derivative]
\lean{FABL.derivative_sq_sub_self_nonneg}
\label{FABL.derivative_sq_sub_self_nonneg}
\leanok
\uses{FABL.BooleanFunction, FABL.BooleanFunction.toReal, FABL.discreteDerivative, FABL.discreteDerivative_apply, FABL.setCoordinate}
\difficulty{4}
Let $n$ be a nonnegative integer, let $f$ be a sign-valued function on $\{-1,1\}^n$, and
choose a coordinate $i$ and an input $x\in\{-1,1\}^n$. The square of the discrete
derivative of the real-valued version of $f$ at $x$, minus that derivative, is
nonnegative.
\end{theorem}

\begin{theorem}[Nonnegativity of the derivative square plus derivative]
\lean{FABL.derivative_sq_add_self_nonneg}
\label{FABL.derivative_sq_add_self_nonneg}
\leanok
\uses{FABL.BooleanFunction, FABL.BooleanFunction.toReal, FABL.discreteDerivative, FABL.discreteDerivative_apply, FABL.setCoordinate}
\difficulty{4}
Let $n$ be a nonnegative integer, let $f$ be a sign-valued function on $\{-1,1\}^n$, and
choose a coordinate $i$ and an input $x\in\{-1,1\}^n$. The square of the discrete
derivative of the real-valued version of $f$ at $x$, plus that derivative, is
nonnegative.
\end{theorem}

\begin{theorem}[Equality characterizes coordinate unateness]
\lean{FABL.abs_fourierCoeff_singleton_eq_influence_iff_isUnateInCoordinate}
\label{FABL.abs_fourierCoeff_singleton_eq_influence_iff_isUnateInCoordinate}
\leanok
\uses{FABL.BooleanFunction, FABL.BooleanFunction.toReal, FABL.IsUnateInCoordinate, FABL.derivative_sq_add_self_nonneg, FABL.derivative_sq_sub_self_nonneg, FABL.discreteDerivative, FABL.discreteDerivative_apply, FABL.fourierCoeff, FABL.influence, FABL.mean, FABL.mean_discreteDerivative_eq_fourierCoeff_singleton, FABL.setCoordinate}
\difficulty{6}
Let $n$ be a nonnegative integer, let $f\colon\{\pm1\}^n\to\{\pm1\}$ be a Boolean
function, and let $i$ be a coordinate. The absolute value of the singleton Fourier
coefficient in coordinate $i$ equals the influence of $f$ in that coordinate if and only
if $f$ is monotone or antimonotone in coordinate $i$.
\end{theorem}

\begin{theorem}[Singleton Fourier coefficients and majority influence]
\lean{FABL.sum_abs_fourierCoeff_singleton_le_totalInfluence_majority}
\proofsource{Analysis-of-Boolean-Functions-by-Ryan-ODonnell}{
  pdf-ffbe9c7b3306-p051-b001,
  pdf-ffbe9c7b3306-p061-b008
}
\label{FABL.sum_abs_fourierCoeff_singleton_le_totalInfluence_majority}
\leanok
\uses{FABL.BooleanFunction, FABL.BooleanFunction.toReal, FABL.Sign, FABL.fourierCoeff, FABL.majority, FABL.majority_monotone, FABL.monomial, FABL.reorientInputEquiv, FABL.signValue, FABL.signValue_thresholdSign, FABL.sum_fourierCoeff_singleton_le_majority, FABL.thresholdSign, FABL.totalInfluence, FABL.totalInfluence_eq_sum_fourierCoeff_singleton_of_monotone}
\difficulty{4}
Let $n$ be a natural number and let $f$ be any Boolean function on the $n$-dimensional
sign cube. The sum, over coordinates $i$, of the absolute values of the singleton
Fourier coefficients of the real-valued version of $f$ is at most the total influence of
the majority function on the same cube.
\end{theorem}

\begin{theorem}[Total influence of unate functions]
\lean{FABL.totalInfluence_toReal_le_majority_of_unate}
\label{FABL.totalInfluence_toReal_le_majority_of_unate}
\leanok
\uses{FABL.BooleanFunction, FABL.BooleanFunction.toReal, FABL.IsUnate, FABL.abs_fourierCoeff_singleton_eq_influence_iff_isUnateInCoordinate, FABL.fourierCoeff, FABL.influence, FABL.majority, FABL.sum_abs_fourierCoeff_singleton_le_totalInfluence_majority, FABL.totalInfluence}
\difficulty{3}
Let $n$ be a nonnegative integer, and let $f:\\{-1,1\\}^n\\to\\{-1,1\\}$ be unate,
meaning that in each coordinate it is either nondecreasing or nonincreasing when that
coordinate changes from $-1$ to $1$. Then the total influence of $f$ is at most the
total influence of majority on $n$ signs.
\end{theorem}

\begin{theorem}[Influence bound for unate functions]
\lean{FABL.totalInfluence_toReal_le_sqrt_card_of_unate}
\proofsource{Analysis-of-Boolean-Functions-by-Ryan-ODonnell}{pdf-ffbe9c7b3306-p063-b002}
\label{FABL.totalInfluence_toReal_le_sqrt_card_of_unate}
\leanok
\uses{FABL.BooleanFunction, FABL.BooleanFunction.toReal, FABL.IsUnate, FABL.abs_fourierCoeff_singleton_eq_influence_iff_isUnateInCoordinate, FABL.fourierCoeff, FABL.fourierWeight, FABL.fourierWeightAtLevel, FABL.fourierWeightAtLevel_one_eq_sum_singleton, FABL.influence, FABL.sum_sq_fourierCoeff_eq_one, FABL.totalInfluence}
\difficulty{4}
Let $n$ be a nonnegative integer, and let $f:\\{-1,1\\}^n\\to\\{-1,1\\}$ be unate in
every coordinate, meaning that each coordinate is either nondecreasing or nonincreasing.
The total influence of $f$ is at most $\\sqrt{n}$.
\end{theorem}

\begin{theorem}[Unateness of linear threshold functions]
\lean{FABL.isUnate_of_isLinearThreshold}
\proofsource{Analysis-of-Boolean-Functions-by-Ryan-ODonnell}{
  pdf-ffbe9c7b3306-p061-b009,
  pdf-ffbe9c7b3306-p132-b002
}
\label{FABL.isUnate_of_isLinearThreshold}
\leanok
\uses{FABL.BooleanFunction, FABL.IsLinearThreshold, FABL.IsUnate, FABL.monotone_thresholdSign, FABL.setCoordinate, FABL.setCoordinate_apply_of_ne, FABL.setCoordinate_apply_self, FABL.signValue, FABL.signValue_neg_one, FABL.signValue_one}
\difficulty{4}
Let $f$ be a sign-valued function on the $n$-dimensional sign cube, and suppose there
are real weights $a_0,a_1,\ldots,a_n$ such that $f(x)$ is the sign determined by whether
$a_0+\sum_{i=1}^n a_i x_i$ is nonnegative. Then, in every coordinate, $f$ is either
monotone or antimonotone: changing that coordinate from $-1$ to $1$ never decreases the
function, or changing it from $1$ to $-1$ never decreases the function.
\end{theorem}
